% Created 2026-09-29 Tue 00:55
% Intended LaTeX compiler: pdflatex
\documentclass[11pt]{book}
\usepackage[utf8]{inputenc}
\usepackage[T1]{fontenc}
\usepackage{graphicx}
\usepackage{longtable}
\usepackage{wrapfig}
\usepackage{rotating}
\usepackage[normalem]{ulem}
\usepackage{amsmath}
\usepackage{amssymb}
\usepackage{capt-of}
\usepackage{hyperref}
\author{The Kalem Authors}
\date{2026-09-29}
\title{Writing Books in Plain Text\\\medskip
\large A sample chapter for Kalem}
\hypersetup{
 pdfauthor={The Kalem Authors},
 pdftitle={Writing Books in Plain Text},
 pdfkeywords={},
 pdfsubject={},
 pdfcreator={Emacs 29.3 (Org mode 9.7.11)}, 
 pdflang={English}}
\begin{document}

\maketitle
\setcounter{tocdepth}{2}
\tableofcontents

This book is written in Org and edited with \emph{Kalem}. It keeps each
chapter in a file of its own; this file pulls them together with
\texttt{\#+INCLUDE} and holds the bibliography.
\part{Foundations}
\label{sec:org89a03a2}
\chapter{Why plain text}
\label{sec:orgb5ff499}
Books have been typeset from plain text for decades (Knuth, Donald E., 1984).
Org keeps the structure of a document in its text: headings, lists,
tables and formulas are all written as characters on a line, so the file
outlives every program that edits it.\footnote{Plain text also diffs well, so a book can be reviewed like code.}
\section{A pipeline}
\label{sec:org6c7538d}
The text goes through one exporter for each format, as
\ref{fig:org0711760} shows: HTML for the web, \LaTeX{} for print, and Word
through pandoc for colleagues who ask for it (Dominik, Carsten, 2010).

\begin{figure}[htbp]
\centering
\includegraphics[width=0.6\textwidth]{figures/pipeline.png}
\caption{\label{fig:org0711760}From one Org file to several formats.}
\end{figure}

The code that builds the book can live in the book itself, as in
literate programming (Schulte, Eric and Davison, Dan and Dye, Thomas and Dominik, Carsten, 2012):

\begin{verbatim}
kalem export book.org --to html
kalem export book.org --to pdf
kalem export book.org --to docx
\end{verbatim}
\captionof{figure}{\label{lst:orgc95557e}Building every format.}
\chapter{Measurements}
\label{sec:org8ed545c}
\ref{tab:orgf625274} lists how long each export of this sample took, and the
total is computed by the table itself. The energy of a body at rest,
equation \ref{eq:orgc7843a4}, is the usual example of a displayed
formula.\footnote{See (Knuth, Donald E., 1984 p. 168) for how \TeX{} sets displayed
equations.}

\begin{table}[htbp]
\caption{\label{tab:orgf625274}Export times of the sample, in milliseconds.}
\centering
\begin{tabular}{lrrr}
Format & First run & Second run & Mean\\
\hline
HTML & 12 & 10 & 11\\
\LaTeX{} & 15 & 13 & 14\\
Word & 420 & 400 & 410\\
\hline
Total & 447 & 423 & 435\\
\end{tabular}
\end{table}

\begin{equation}
\label{eq:orgc7843a4}
E = m c^2
\end{equation}

Inline mathematics works too: the mean of \(n\) runs is
\(\bar{x} = \frac{1}{n}\sum_{i=1}^{n} x_i\). Back in
\hyperref[sec:orgb5ff499]{the first chapter}, the pipeline of
\ref{sec:org6c7538d} produced these files.
\part*{References}
\label{sec:org44d2ff9}
\noindent
Dominik, Carsten (2010). \emph{The Org Mode 7 Reference Manual}, Network Theory.

\noindent
Knuth, Donald E. (1984). \emph{The {\TeX}book}, Addison-Wesley.

\noindent
Schulte, Eric and Davison, Dan and Dye, Thomas and Dominik, Carsten (2012). \emph{A Multi-Language Computing Environment for Literate Programming and Reproducible Research}.
\end{document}